\subsection{角}

由于域 \(\mathbf{R}\) 是有序的，我们可以定向实数平面 \(\mathbf{R}^2\) ，方法是取 \(e_1 \wedge e_2\) 为正二重向量（\(e_1, e_2\) 是标准基的向量）。于是，在定向的实数平面 \(\mathbf{R}^2\) （以下与 \(\mathbf{C}\) 等同）中，我们可以对任意一对射线定义角 \((\widehat{\Delta_1, \Delta_2})\) （射线对（ \(\Delta_1, \Delta_2\) ）以 \(o\) 为端点）(*)。所有角组成的集合 \(\mathfrak{A}\) 具有一个交换群结构（写成加法），由下式定义

\[
(\widehat {\Delta_ {1} , \Delta_ {3}}) = (\widehat {\Delta_ {1} , \Delta_ {2}}) + (\widehat {\Delta_ {2} , \Delta_ {3}}),
\]

于是特别有 \((\widehat{\Delta_1, \Delta_1}) = 0\) 且 \((\widehat{\Delta_2, \Delta_1}) = -(\widehat{\Delta_1, \Delta_2})\) 。

平角 \(\varpi\) 是 \(\neq 0\) 的解，即方程 \(2\theta = 0\) 在 \(\mathfrak{A}\) 中的解；它是负实半轴与正实半轴所成的角。

若 z 是任一非零复数，z 的辐角（或幅角），记作 \(\operatorname{Am}(z)\) ，是过 z 而以 o 为端点的射线与正实半轴所成的角。映射 \(z \to \operatorname{Am}(z)\) 是乘法群 \(C^{*}\) 到加法群 A 上的一个同态，因而我们有

\[
\operatorname{Am} \left(z z ^ {\prime}\right) = \operatorname{Am} (z) + \operatorname{Am} \left(z ^ {\prime}\right) \quad \text { and } \quad \operatorname{Am} (\bar {z}) = \operatorname{Am} \left(z ^ {- 1}\right) = - \operatorname{Am} (z).
\]

角 \(\delta = \operatorname{Am}(i)\) 称为正直角；它是方程 \(2\theta = \varpi\) 在 A 中的解之一，另一个解是 \(-\delta = \delta + \varpi\) 。

同态 \(z \to \operatorname{Am}(z)\) 限制到子群 \(\mathbf{U}\) （\(\mathbf{C}^*\) 的）上时，是 \(\mathbf{U}\) 的群结构到 \(\mathfrak{A}^{(**)}\) 的群结构上的一个同构；若用这个同构把 U 的拓扑转移到群 A 上，则后一群成为一个紧拓扑群，并且映射 \(z \to \operatorname{Am}(z)\) （\(C^{*}\) 到 A 上的）是拓扑群 \(C^{*}\) 到拓扑群 A 上的一个严格态射。

我们把 \(\theta \to f(\theta)\) 记为 \(\mathfrak{A}\) 到 \(\mathbf{U}\) 上的同构，它即同构 \(z \to \operatorname{Am}(z)\) （\(\mathbf{U}\) 到 \(\mathfrak{A}\) 上的）的逆。按定义，\(\mathcal{R}(f(\theta))\) 记作 \(\cos \theta\) ，称为角 \(\theta\) 的余弦；\(\mathfrak{J}(f(\theta))\) 记作 \(\sin \theta\) ，称为角 \(\theta\) 的正弦。这些函数在拓扑群 \(\mathfrak{A}\) 上连续，并且满足下列关系式（同处所引），它们是上述定义的直接推论：

\[
\begin{array}{r l} \cos o = 1, & \sin o = 0, \quad \cos \varpi = - 1, \quad \sin \varpi = 0, \\ \cos (- \theta) = \cos \theta , & \sin (- \theta) = - \sin \theta , \\ \cos (\theta + \theta^ {\prime}) = \cos \theta \cos \theta^ {\prime} - \sin \theta \sin \theta^ {\prime}, \\ \sin (\theta + \theta^ {\prime}) = \sin \theta \cos \theta^ {\prime} + \sin \theta^ {\prime} \cos \theta , \\ \cos^ {2} \theta + \sin^ {2} \theta = 1. \end{array}
\]

按定义，角 \(\theta \in \mathfrak{A}\) 的正切，当 \(\cos \theta \neq 0\) 时定义为 \(\sin \theta / \cos \theta\) （同处所引），记作 \(\tan \theta\) ；它是一个连续函数，由连续性延拓到 \(\tilde{\mathbf{R}}\) （第六章，§ 3，第 4 号），即取值 \(\infty\) 于角 \(\delta\) 与 \(-\delta\) 处。我们有 \(\tan (\theta + \varpi) = \tan \theta\) 。\(\theta\) 的余切，记作 \(\cot \theta\) ，是 \(\tilde{\mathbf{R}}\) 中等于 \(1 / \tan \theta\) 的元素。

注意，若 \(\operatorname{Am}(z) = \theta\) ，我们有 \(z = |z| (\cos \theta + i \sin \theta)\) ；这个表达式称为复数 \(z \neq 0\) 的三角形式。

